% ECONOMICS_PROBLEM: {"id": "D62-140", "title": "Data externalities", "jel_code": "D62", "source_id": "OP-0407"}
% !TeX program = xelatex
\documentclass[11pt,a4paper]{article}
\usepackage[margin=23mm]{geometry}
\usepackage{fontspec}
\setmainfont{Latin Modern Roman}
\usepackage{amsmath,amssymb,mathtools}
\usepackage{xcolor,enumitem,booktabs,longtable,array,xurl,hyperref}
\definecolor{muted}{HTML}{53636C}
\hypersetup{colorlinks=true,urlcolor=blue,linkcolor=blue}
\setlength{\parindent}{0pt}
\setlength{\parskip}{5pt}
\newcommand{\R}{\mathbb R}
\newcommand{\E}{\mathbb E}
\newcommand{\N}{\mathbb N}
\newcommand{\ind}{\mathbf 1}
\newcommand{\Var}{\operatorname{Var}}
\newcommand{\Cov}{\operatorname{Cov}}
\newcommand{\supp}{\operatorname{supp}}
\DeclareMathOperator*{\argmin}{arg\,min}
\DeclareMathOperator*{\argmax}{arg\,max}
\newcommand{\entrymeta}[3]{{\sffamily\small\color{muted}\textbf{\texttt{#1}}\quad|\quad #2\par #3\par}}
\newcommand{\Problemref}[1]{\texttt{#1}}
\newcommand{\JELref}[2]{\texttt{#2} (JEL \texttt{#1})}
\begin{document}
\section*{Data externalities}
\textbf{Problem ID:} \texttt{D62-140}\quad \textbf{JEL:} \texttt{D62}\par
% BEGIN ECONOMICS STATEMENT
\label{problem:OP-0407}
{\sffamily\small\color{muted}Original subject group: Digital, platform and data economics\par}
\label{definitions:problem:30007689}
\textbf{Audit classification:} Background; proposed extension. The official chapter and closing research topics were verified through primary indexing; this rights-design mechanism is a new specification. Its fully known finite benchmark is computable by enumeration.
\subsection*{Definitions and precise research target}
Consider a finite population of \(n\) consumers whose private characteristics are jointly distributed according to a known prior \(P\). Consumer \(i\) chooses a disclosure indicator \(d_i\in\{0,1\}\); the vector \(d\) determines signals available to a data intermediary. Correlation means that signals about one consumer may change posterior beliefs about another. A downstream firm uses those posteriors to choose products or prices. Let \(v_i(d)\) be consumer \(i\)'s expected service benefit minus payments and privacy losses, and let \(\pi(d)\) be firms' expected profit net of real data-processing costs. Define total welfare \(W(d)=\sum_{i=1}^n v_i(d)+\pi(d)\). A rights mechanism \(m\) specifies consent, compensation and allowable downstream uses; \(\mathcal M\) is a stated feasible class with voluntary participation and truthful reporting. Under \(m\), disclosure choices form equilibrium set \(\mathcal E(m)\). Seek mechanisms maximizing a declared selection-weighted expectation of \(W(d)\) over \(d\in\mathcal E(m)\), subject to each nonconsenting person's specified privacy-risk bound. Compare individual consent, collective representation and limits on inference about nonparticipants. Quantify informational spillovers separately from transfers. The specification requires a clear privacy-loss metric and equilibrium rule; it is an original design proposal motivated by data externalities, without a verified canonical open formulation.

\textbf{Definitions, assumptions and mathematical qualifications.} The prior covers latent types and all signals jointly. A data-use rule maps revealed signals to downstream actions; a privacy-risk functional \(\rho_i(d,a)\) must be supplied, for example expected monetary harm rather than unspecified privacy. A mechanism defines transfer funding, consent contingencies and agents information; truthful reporting and voluntary participation are constraints for each type or in expectation, as explicitly chosen. Fix \(\kappa_m\), a probability distribution on a nonempty equilibrium set, and define \(\mathcal W(m)=\int W(d)d\kappa_m(d)\). Constraint \(\rho_i\leq\bar\rho_i\) applies to each nonconsenting individual. A finite mechanism menu has a maximizer only when the feasible subset is nonempty; otherwise use a supremum with existence conditions.
\subsection*{Variants and assumption changes}
\begin{enumerate}
\item \textbf{Editorial, independent types.} Set cross-person statistical dependence to zero while retaining any service or price spillovers; identify which externalities disappear and which remain.
\item \textbf{Editorial, correlated types.} Permit inference about nonparticipants, compare individual and collective consent mechanisms using a supplied loss metric, and optimize robust welfare over admissible signal and harm distributions.
\end{enumerate}
\subsection*{References and source audit}
\textbf{Checked source.} The Economics of Privacy, Chapter 3 (2024), pp. 73--96; indexed closing passage p. 93. The indexed conclusion names data combination and regulatory evaluation as further work. Primary chapter metadata and indexed passage checked; full PDF returned HTTP 403. The rights-mechanism question is editorial. \url{https://www.nber.org/system/files/chapters/c14782/c14782.pdf}.
\begin{enumerate}\raggedright
\item\hypertarget{definitions:source:30007689:1}{} Alessandro Bonatti. 2024. \emph{The Platform Dimension of Digital Privacy}. The Economics of Privacy, Chapter 3 (2024), pp. 73--96; indexed closing passage p. 93.
\url{https://www.nber.org/system/files/chapters/c14782/c14782.pdf}.
\end{enumerate}

\subsection*{Common conventions: Source mathematical conventions}

These conventions apply to each entry unless explicitly replaced.

\textbf{Models and identification.} A primitive model \(m\) specifies all probability laws, feasible choices, information, equilibrium and measurement rules used by its target. The admissible class \(\mathcal M\) is fixed independently of outcomes. If \(P_O(m)\) is its observable probability law and \(\tau(m)\) the target, its identified set at observable law \(P\) is
\[
 \mathcal I_\tau(P)=\{\tau(m):m\in\mathcal M,\ P_O(m)=P\}.
\]
Point identification means that this set is a singleton. An empty set signals incompatibility of data and assumptions. Coordinate bounds need not describe the joint set. Bounds are sharp only if every retained target is attainable or arbitrarily approachable by compatible models. Finite bounds require explicit restrictions on unobserved quantities; unrestricted confounding, hidden wealth, extrapolation or arbitrary equilibrium selection can yield uninformative sets.

\textbf{Causal interventions.} A potential outcome \(Y_i(p)\) is the outcome under a complete policy regime \(p\), including timing, financing, information and market spillovers. The target population and units are fixed before treatment. With interference, use the complete assignment vector or a justified exposure mapping; individual no-interference notation alone cannot identify an economy-wide intervention. A difference of mean potential outcomes does not require the cross-world joint distribution. Conditioning on post-treatment migration, survival, enrollment or employment changes the estimand and needs separate justification.

\textbf{Statistical statements.} An estimator, confidence set, forecast or test specifies a sampling law, model class and loss/coverage criterion. Uniform \((1-\alpha)\)-coverage means \(\inf_{m\in\mathcal M}\Pr_m\{\tau(m)\in C_n\}\geq1-\alpha\), or an explicitly asymptotic version. The whole parameter space trivially has coverage, so informativeness requires a stated width, risk or power objective. A forecasting improvement is relative to a named benchmark, information set and evaluation population.

\textbf{Equilibrium and optimization.} An equilibrium comprises feasible plans, optimality, beliefs where needed, budget constraints and all market/resource clearing. The primitives or a stated benchmark define each correspondence \(\mathcal E(p;m)\). Nonemptiness is assumed or proved, never inferred just from compact policy sets. A selected-equilibrium target uses a declared selection rule; a robust target quantifies over all permitted equilibria. Use suprema or infima unless attainment is established. An \(\varepsilon\)-optimal policy is feasible and within \(\varepsilon\) of the finite optimum value. Report an empty feasible set separately.

\textbf{Welfare and accounting.} Normative utility, interpersonal normalization, social weights, numeraire, discounting and the use of public receipts are primitives. Behavioral utility need not coincide with the welfare criterion. Household welfare includes ownership income and rents once; adding profits again to compensating variation generally double-counts them. A monetary equivalent is defined by an explicit transfer and price convention, with monotonicity and range conditions. Average effects, mechanism contributions, transfers and welfare losses are different targets. Infinite sums require convergence and dynamic feasible plans require terminal or no-Ponzi conditions.

\textbf{Source versions.} A working-paper date, revision date and journal publication year are distinguished. A DOI for a journal version is not automatically the DOI of an earlier manuscript. Locators refer to the stated version and printed pages where specified. A metadata-only check does not verify a claim about a full-text conclusion. Every entry's source subsection narrows the provenance claim accordingly.
% END ECONOMICS STATEMENT
\end{document}
